\section*{Mentoring Plan}
\addcontentsline{toc}{section}{\protect\numberline{}Mentoring Plan}

The proposed research requires the broad background of a research associate with expertise in mathematical statistics as well as experience in machine learning and programming. These include the development of software, strong foundation knowledge in the STEM fields, research methodology, experimental design and analysis. 
% Researchers will be encouraged to work and interact with other research groups to gain a broad understanding of the CS and broader stem fields including psychology, learning sciences, and more. The interdisciplinary research experience will provide the student with a unique research educational development.

The PI will provide support including some areas listed as follows:

\noindent\textbf{Weekly Discussions: } Students will meet with the PI at least once a week for discussions. These include research progress, brainstorming, goal setting, and professional counseling activities. Additionally, students will be expected to present their findings on a regular basis in front of a large audience in UTK's Data Science and Statistics Seminar. This will develop their presentation skills as well. 

\noindent\textbf{Publications and Presentation: } Students will be the primary lead authors on Thrust I. The PI's role will be to support the student and make sure that the projects are completed timely and that the work being conducted is of high quality. For Thrust II and III, the PI will first work with the students to layout the theoretical development. The students will be expected to learn from this process and meet with the PI and collaborators regularly to ensure that their project is of the highest quality.

\noindent\textbf{Peer Review: } Students will be trained to provide proper peer reviews. Students are expected to present within the department regularly and seek out the feedback of their peers. As the projects shape up, students will be encouraged to attend conferences and present their work. They will also be encouraged to register as peer review for conferences. 

\noindent\textbf{Technical Development: } Through working in the proposed topics, students will develop their technical skills in Mathematical Statistics, Differential Equations, Machine learning and applications. 
% They will also learn how to conduct interdisciplinary research and propose creative use cases for human-ai collaborative scenarios. Proper and ethical experimental research methods (e.g. user studies) will also be learned.

\noindent\textbf{Connections to Industry: } Students will have the opportunity to work with long time industry collaborators of the PI as well as establish connections through internships. It is expected that students seek out internships during summers especially in their 3rd and 4th years.

\noindent\textbf{Orientation to the Project:} 
The student researcher will have an initial meeting with the PI, to whom the student researcher will report. Expectations for their work on the project will be discussed and put into writing. These include the amount of supervisory responsibility, independent research projects, expected publications, and grant writing expectations. 
